\documentclass[year1]{../idcore-ar}
% Note: If compiling in the same directory as idcore-ar.cls, use:
% \documentclass[year1]{idcore-ar}
% To suppress guidance callout boxes in your final submission, add the option:
% \documentclass[year1, hideguidance]{../idcore-ar}

% -----------------------------------------------------------------------------
% Report & Student Metadata
% -----------------------------------------------------------------------------
\author{Jane Doe}
\studentnumber{s1234567}
\projecttitle{Next-Generation Floating Offshore Wind Turbine Dynamics in Complex Wave Climates}

% Supervisory Team
\leadsupervisor{Prof. David Ingram (University of Edinburgh)}
\cosupervisorone{Prof. Julia Race (University of Strathclyde)}
\cosupervisortwo{Dr. Philipp Thies (University of Exeter)}
\independentreviewer{Prof. Ian Masters (Swansea University)}
\industrialsupervisor{Dr. Robert Miller}
\company{Renewable Dynamics Ltd}

% Logo Path
\logopath{../../assets/idcore_logo.png}

% Year 1 Specific: Research Ethics & Integrity Assessment
\ethicsrequired{No} % Set to 'Yes' or 'No'
\ethicsattached{No} % Set to 'Yes' or 'No'

% Document Control Revisions (Add multiple revision lines using \addrevision{version}{date}{amendment})
% To start fresh, call \clearrevisions before adding your revisions.
\clearrevisions
\addrevision{1.0}{15/07/2026}{Sent to supervisors for review}
\addrevision{2.0}{20/08/2026}{Updated following review discussion}

% Signatures & Approvals (Text or signature placeholder)
\studentsignature{J. Doe}
\signaturedate{\today}
\industrialsupervisorsignature{R. Miller}
\industrialsupervisordate{\today}
\horiname{Prof Alasdair McDonald}
\horisignature{A. McDonald}
\horidate{\today}

% -----------------------------------------------------------------------------
% Document Body
% -----------------------------------------------------------------------------
\begin{document}

% Generate official Year 1 Cover Sheet
\makecoverpage

% =============================================================================
% 1. Literature Review
% =============================================================================
\section{Literature Review}

\begin{guidance}[Page Limit: 3 pages only | Minimum 11pt font]
\textbf{Purpose:} Critical review of current literature and state of the art.\\
\textbf{Guidance Points:}
\begin{itemize}\setlength{\itemsep}{1pt}
  \item Write for a technical, but non-specialist reader.
  \item Establish where current knowledge ends and where your research question begins.
  \item Synthesise existing academic and industrial methodologies, highlighting limitations and research gaps that your EngD project will address.
  \item Maintain consistent Harvard referencing throughout (e.g., \citet{smith2024tidal}; \citet{macdonald2022offshore}).
\end{itemize}
\end{guidance}

[Enter your Literature Review here. Maintain maximum 3 pages.]


% =============================================================================
% 2. Project Plan
% =============================================================================
\section{Project Plan}

\begin{guidance}[Page Limit: 3 pages only | Minimum 11pt font]
\textbf{Purpose:} Define your research strategy for the remaining 3 years of the EngD.\\
\textbf{Key Areas to Cover:}
\begin{enumerate}\setlength{\itemsep}{1pt}
  \item \textbf{What do you plan to do and why?} Anticipated contribution to knowledge, novelty, importance, deliverables (kit, code, or theory), and clear statement of overall Aims.
  \item \textbf{How will you achieve this?} Specific stages and precise list of measurable objectives.
  \item \textbf{What resources are required?} Lab facilities (e.g., FloWave, FastBlade), HPC access, equipment, data sets, software licences, travel funds, external partner reliance.
  \item \textbf{Risks to successful completion:} Realistic evaluation of dependencies, unexpected experimental results, delays in data or materials, and mitigation alternatives.
\end{enumerate}
\end{guidance}

\subsection{Contribution to Knowledge and Project Aims}
[State your anticipated novelty, tangible deliverables, alignment with UK Net Zero targets, and primary Aims.]

\subsection{Methodology and Measurable Objectives}
[Detail the project stages and provide a bulleted list of measurable, verifiable objectives.]

\subsection{Required Resources and Collaborations}
[List required experimental facilities, computing clusters, software licences, and sponsor data access.]

\subsection{Risks to Successful Completion}
[Analyse potential obstacles, external dependencies, and mitigation strategies.]


% =============================================================================
% 3. References
% =============================================================================
\newpage
\section{References}

\begin{guidance}[Page Limit: 1 page only]
All IDCORE students must use the \textbf{Harvard Referencing Style} (author-date format).
Ensure consistent formatting throughout your written work all the way through to your final EngD thesis.
\end{guidance}

\bibliographystyle{abbrvnat}
\bibliography{../references}


% =============================================================================
% 4. Gantt Chart & Work Plan
% =============================================================================
\newpage
\section{Gantt Chart \& Diagrammatic Work Plan}

\begin{guidance}[Page Limit: 1 page only]
Include a diagrammatic work plan (Gantt chart) on this dedicated page:
\begin{itemize}\setlength{\itemsep}{1pt}
  \item Show Year 1 / upcoming 12 months in reasonable detail.
  \item Outline milestones, journal/conference deliverables, experimental test windows, and taught module elements (60 research credits during the research phase).
  \item \textbf{Mandatory Requirement:} Your plans must include an explicit allocation of time for writing up your thesis/portfolio --- \textbf{ideally at least 6 months}.
\end{itemize}
\end{guidance}

\begin{center}
  \vspace{1cm}
  % Insert your Gantt Chart figure here, e.g.:
  % \includegraphics[width=0.95\linewidth]{figures/gantt_chart.pdf}
  \fbox{\parbox[c][12cm][c]{0.9\linewidth}{\centering\large\color{idcoregray}\textbf{[Insert 1-Page Gantt Chart Here]}\\
  \vspace{4mm}
  \small Must include Year 1 detail and at least 6 months reserved for thesis write-up.}}
\end{center}


% =============================================================================
% 5. Project Risk Assessment
% =============================================================================
\newpage
\section{Project Risk Assessment}

\begin{guidance}[Requirement: 5 key project risks]
Identify and evaluate the \textbf{5 key risks} to successfully completing your research project within the prescribed 3-year research window.
\end{guidance}

\begin{table}[htbp]
\centering
\small
\renewcommand{\arraystretch}{1.3}
\begin{tabularx}{\linewidth}{|c|>{\raggedright\arraybackslash}p{4cm}|c|c|>{\raggedright\arraybackslash}X|}
  \hline
  \rowcolor{idcorelight}
  \textbf{\#} & \textbf{Risk Description} & \textbf{Likelihood} & \textbf{Impact} & \textbf{Mitigation Strategy} \\ \hline
  1 & Delay in test tank booking / availability & Medium & High & Engage facility managers early; secure backup slots; begin with numerical modelling. \\ \hline
  2 & Sponsor proprietary sensor data delay & Medium & Medium & Establish fallback data agreements; use public/synthetic benchmarking datasets. \\ \hline
  3 & Numerical simulation solver convergence issues & Low & Medium & Validate against published empirical benchmarks; seek academic specialist input. \\ \hline
  4 & High-performance computing queue bottlenecks & Low & Low & Optimise code locally; utilise partner clusters (e.g., HPC-Wales / Cirrus). \\ \hline
  5 & Unforeseen scope expansion by industrial sponsor & Medium & High & Maintain fortnightly supervisory reviews; enforce supervision agreement scope. \\ \hline
\end{tabularx}
\caption{Key Project Risks and Mitigation Plans.}
\end{table}


% =============================================================================
% 6. Health \& Safety Assessment
% =============================================================================
\section{Health \& Safety Assessment}

\begin{guidance}[Length: Not restricted in length]
Detail all relevant Health \& Safety considerations for your project:
\begin{itemize}\setlength{\itemsep}{1pt}
  \item Company-specific workplace, lab, and offshore site visit policies.
  \item Laboratory or test tank testing protocols (electrical safety, high pressure, heavy lifting, flume operations).
  \item Lone working and fieldwork risk assessments.
\end{itemize}
\end{guidance}

[Detail your project's Health \& Safety assessment here.]


% =============================================================================
% 7. Responsible Research and Innovation (RRI) Report
% =============================================================================
\section{Responsible Research and Innovation (RRI) Report}

\begin{guidance}[Length: Not restricted in length]
The RRI report reflects on the broader societal, environmental, ethical, and economic implications of your research.\\
\textbf{Key Guidance:}
\begin{itemize}\setlength{\itemsep}{1pt}
  \item Address the AREA framework: \textbf{Anticipate} outcomes, \textbf{Reflect} on motivations and assumptions, \textbf{Engage} with stakeholders and industry, and \textbf{Act} accordingly.
  \item If the RRI or Project Plan indicate ethical issues with your research, you \textbf{must} complete an additional Research Ethics and Integrity form (see Cover Sheet and School requirements at \url{http://www.eng.ed.ac.uk/research/ethics/}).
\end{itemize}
\end{guidance}

[Detail your RRI reflection and assessment here.]


% =============================================================================
% 8. Trusted Research & Innovation (TR&I) and Data Management Plan
% =============================================================================
\newpage
\section{Trusted Research \& Innovation (TR\&I) and Data Management Plan}

\begin{guidance}[Requirement: Initial Data & Security Governance | Lighter Touch]
The IDCORE Industrial Advisory Board (IAB) requires all research engineers to establish strong data governance, IP protection, and research security protocols from project kick-off.\\
\textbf{Key Planning Areas:}
\begin{enumerate}\setlength{\itemsep}{1pt}
  \item \textbf{Data Management Plan (DMP):} Types and estimated volumes of data (experimental, sensor logs, numerical simulations, CAD models); file formats; backup redundancy and storage location (company secure drives, university cloud).
  \item \textbf{Sponsor Intellectual Property (IP) Baseline:} Agreed treatment of background IP, ownership of anticipated foreground IP, non-disclosure agreements (NDAs), and boundaries of commercial confidentiality.
  \item \textbf{Trusted Research Screening:} Initial screening under UK Government / NPSA Trusted Research principles (\url{https://www.npsa.gov.uk/specialised-guidance/trusted-research}) regarding international collaborations, sensitive software, or external data sharing.
\end{enumerate}
\end{guidance}

\subsection{Data Management Plan (DMP) and Storage Architecture}
[Describe datasets to be generated, storage locations, backup mechanisms, and data organisation.]

\subsection{Sponsor Intellectual Property (IP) and Confidentiality Baseline}
[Detail agreement with industrial sponsor regarding background IP, foreground IP, and non-disclosure boundaries.]

\subsection{Trusted Research \& Innovation (TR\&I) Preliminary Screening}
[Summarise any international research security considerations, sensitive dual-use technology screening, or data access controls.]

\end{document}
